この章では、本システムにて使用したプログラミング言語の特徴・役割について説明する。
\section{HTML}
HTML（Hyper Text Markup Language）とは、Webページの作成に使われるマークアップ言語であり、文章の構造や画像、リンクなどを定義することで、Webページの骨組みを記述するために用いられる基本的な言語である。
\par
HTMLは主に、head要素とbody要素の二つから構成される。head要素には、文書に関するメタ情報が記述され、通常は画面上に表示されない。ただし、title要素に記述された内容は、Webブラウザのタイトルバーやタブのページタイトルとして表示される。
また、head要素内では、CSSファイルやJavaScriptファイルなどの外部ファイルを読み込むことが可能である。
body要素には、Webブラウザ上に表示されるすべての内容が記述される。テキストや画像、ボタンなどの要素が含まれ、JavaScriptファイルの読み込みやJavaScriptプログラムの記述もおこなうことができる。
\section{JavaScript}